\documentclass[11pt]{article}
\usepackage[margin=1in]{geometry}
\usepackage{amsmath,amssymb}
\usepackage{graphicx}
\usepackage{hyperref}
\usepackage{authblk}
\usepackage{natbib}

\title{A Color-Dependent Structure in SPARC Rotation-Curve Residuals}
\author{Antoine L. Shephard (Echo Mirrowen)}
\affil{Independent Researcher}
\date{\today}

\begin{document}
\maketitle

\begin{abstract}
Using the public SPARC rotation-curve mass-model decompositions, we define an ``apparent dark residual''
$f_{\rm DM}(r) = 1 - V_{\rm bar}^2(r) / V_{\rm obs}^2(r)$ with
$V_{\rm bar}^2 = V_{\rm gas}^2 + V_{\rm disk}^2 + V_{\rm bulge}^2$.
We summarize $f_{\rm DM}$ at multiple radii and test for correlations with WISE-based stellar mass and the optical--IR color $g-W1$.
Across the SPARC--WISE overlap sample, $f_{\rm DM}$ exhibits a strong monotonic dependence on $g-W1$ that strengthens when restricting to
high-quality rotation curves and disk-dominated systems, and persists at standardized radii ($2R_{\rm eff}$ and $2.2R_d$).
We provide the merged dataset and scripts sufficient to reproduce all figures and statistics.
\end{abstract}

\section{Data}
We use the SPARC sample and associated rotation-curve mass models \citep{Lelli2016}.
For each galaxy we use the SPARC \texttt{rotmod} decomposition (rotation curve and baryonic contributions) to compute $f_{\rm DM}(r)$.
We cross-match galaxies with WISE-based stellar masses and $g-W1$ color from \citep{Duey2025}.
Rotation-curve quality flags and structural radii ($R_{\rm eff}$, $R_d$) are taken from SPARC Table~1 \citep{Lelli2016}.

\section{Residual definition and summary metrics}
For each radius point in the SPARC \texttt{rotmod} files we compute
\begin{equation}
f_{\rm DM}(r) \equiv 1 - \frac{V_{\rm bar}^2(r)}{V_{\rm obs}^2(r)},\qquad
V_{\rm bar}^2(r)=V_{\rm gas}^2(r)+V_{\rm disk}^2(r)+V_{\rm bulge}^2(r).
\end{equation}
We report four summary metrics:
(i) the mean of $f_{\rm DM}$ over the outer 25\% in radius ($\langle f_{\rm DM}\rangle_{\rm outer}$),
(ii) the value at the outermost measured point ($f_{\rm DM, last}$),
(iii) $f_{\rm DM}(2R_{\rm eff})$, and (iv) $f_{\rm DM}(2.2R_d)$ (linear interpolation within each galaxy's measured range).

To reduce morphology confounding, we define a disk-dominated subset using the outer bulge contribution fraction
$\langle V_{\rm bulge}^2 / V_{\rm bar}^2\rangle_{\rm outer} \le 0.05$.
We also test a high-quality subset using the SPARC quality flag $Q=1$.

\section{Results}
Figure~\ref{fig:main} shows the relationship between $\langle f_{\rm DM}\rangle_{\rm outer}$ and $g-W1$ for the full matched sample and for the disk-dominated, $Q=1$ subset.
The monotonic correlation strengthens under these restrictions and remains significant for alternative residual summaries, including standardized radii.
A brief statistical summary is provided in the accompanying \texttt{CSV} artifact and can be reproduced by the provided scripts.

\begin{figure}[t]
\centering
\includegraphics[width=0.48\linewidth]{../../docs/observations/figures/fdm_outer_vs_gw1_all.png}\hfill
\includegraphics[width=0.48\linewidth]{../../docs/observations/figures/fdm_outer_vs_gw1_disk_Q1.png}
\caption{Left: SPARC--WISE matched sample, outer-mean residual $\langle f_{\rm DM}\rangle_{\rm outer}$ versus $g-W1$.
Right: disk-dominated and high-quality subset (outer bulge fraction $\le 0.05$ and SPARC $Q=1$).}
\label{fig:main}
\end{figure}

\section{Discussion}
The observed color dependence implies that the Newtonian mismatch between baryonic mass models and observed rotation curves is structured with
stellar-population / formation proxies (here, $g-W1$) and not purely random scatter.
Because $g-W1$ correlates with stellar mass and morphology, the results should be interpreted as a phenomenological constraint on rotation-curve residual structure.
Future work can extend this analysis with explicit morphology controls and alternative stellar-population indicators.

\section{Reproducibility}
We include (i) a merged dataset (\texttt{docs/observations/data/sparc\_wise\_fdm\_with\_Q\_and\_radii.csv})
containing all derived residual metrics and matched properties, and (ii) repository notes sufficient to reconstruct the analysis pipeline.
All plots in this note can be regenerated from SPARC public data products \citep{Lelli2016} and the WISE stellar-mass tables \citep{Duey2025}.

\bibliographystyle{plainnat}
\bibliography{references}

\end{document}
